%% kapa: exam=1 round=22 year=2011 period=1 form=A subject=economics subject_ko=경제원론 seq=5 number=45
%% source: Q-Net 감정평가사 기출문제 (cst00309) artlSeq=5036790 file=extracted/1008/5036790/2011년 제22회 감정평가사 1차시험 기출문제/1교시.hwp
%% conversion: kordoc 4.24.51 -> tools/gen1.py (auto); review=proofread-sonnet; flags=none
\begin{question}{45}
두 재화 X와 Y의 가격이 제1기에 P\textsubscript{X}=10, P\textsubscript{Y}=40이었으며, 甲은 재화소비조합점 $(x,~y)$=(60, 20)을 선택하였다. 현시선호이론에 관한 다음 설명 중 옳은 것만을 모두 고른 것은? (단, $x,~y$는 각각 X재와 Y재의 소비량, P\textsubscript{X}와 P\textsubscript{Y}는 각각 X재와 Y재의 가격)
\begin{framedtext}
ㄱ. 제2기에 가격이 P\textsubscript{X}=20, P\textsubscript{Y}=30으로 변화했을 때, 甲이 재화소비조합점 (65, 15)를 선택했다면 甲의 선택은 약공리를 위배하지 않는다. \par
ㄴ. 제2기에 가격이 P\textsubscript{X}=20, P\textsubscript{Y}=20으로 변화했을 때, 甲이 재화소비조합점 (50, 30)을 선택했다면 甲의 선택은 약공리를 위배한다. \par
ㄷ. 강공리가 성립하면 약공리는 항상 성립한다.
\end{framedtext}
\begin{choices}
  \item ㄱ
  \item ㄱ, ㄴ
  \item ㄴ, ㄷ
  \item ㄱ, ㄷ
  \item ㄱ, ㄴ, ㄷ
\end{choices}
\end{question}
\answer{4}
